\begin{abstract}
Self-evolving agents convert execution feedback into persistent skills, but most studies emphasize which trajectories reach the updater rather than whether the updater reliably materializes a useful persistent state. We study a narrower failure mode: byte-identical learner-visible evidence can fail to regenerate the same behaviorally useful skill. A controlled DeepSeek-V4-Pro study of 48 matched evidence pairs, 96 learned states, and 1,728 held-out evaluations first shows that replacing winner evidence with a mixed rejected-witness view has a small heterogeneous mean effect of 0.023, with a frozen confidence interval crossing zero. We then examine one outcome-selected development stream. A historically strong First-Fail state remains directionally useful across frozen-state actor remeasurements, whereas two fresh updater states produced from reconstructed byte-identical First-Fail evidence do not reproduce the historical advantage. These observations rule out evidence disappearance and weaken a pure one-rollout actor-luck explanation, while remaining consistent with a persistent-state generation bottleneck; they do not identify a population variance component. We therefore treat state generation as an experimental variable and define a prospective typed compiler intervention that maps the same trajectory+score package to a finite repair representation and canonical skill state. A fresh Evidence$\times$Generator bridge independently tests the matched-evidence generator effect, trajectory-conditioned content against score-only and scope-matched generic controls, and any secondary First-Fail evidence effect. The present paper establishes local state-regeneration instability and a falsifiable generator hypothesis; it does not yet claim that the compiler improves utility or that rejected failures are uniquely beneficial.
\end{abstract}
